\subsection{Affine points as valued residue fields}
\label{subsec:valued-spectrum-points}

These are actual primes and actual valuations on their residue fields,
over any commutative ring. The construction describes the affine point
space and its evaluation topology. It does not identify a general scheme's
analytification or construct an analytic structure sheaf.

\begin{theorem}[Valued residue-field points]
\label{spectrum:valued-point-equivalence}
\lean{EnvelopeProperty.Berkovich.ValuedPoint,
EnvelopeProperty.Berkovich.ValuedPoint.eval,
EnvelopeProperty.Berkovich.ValuedPoint.toSpectrum,
EnvelopeProperty.Berkovich.ValuedPoint.toSpectrum_val_apply,
EnvelopeProperty.Berkovich.ValuedPoint.eval_eq_zero_iff,
EnvelopeProperty.Berkovich.ValuedPoint.support_toSpectrum,
EnvelopeProperty.Berkovich.ValuedPoint.ofSpectrum,
EnvelopeProperty.Berkovich.ValuedPoint.ofSpectrum_prime,
EnvelopeProperty.Berkovich.ValuedPoint.eval_ofSpectrum,
EnvelopeProperty.Berkovich.ValuedPoint.toSpectrum_ofSpectrum,
EnvelopeProperty.Berkovich.ValuedPoint.toSpectrum_injective,
EnvelopeProperty.Berkovich.ValuedPoint.ofSpectrum_toSpectrum,
EnvelopeProperty.Berkovich.ValuedPoint.spectrumEquiv}
\leanok
\uses{spectrum:support,spectrum:residue-field}
For a commutative ring $A$, let a valued point be a pair consisting of
$p\in\operatorname{Spec}(A)$ and a real multiplicative ultrametric valuation
$\nu:\kappa(p)\to\mathbb R_{\ge0}$. Evaluation
$|a|=\nu(a\bmod p)$ gives an actual seminorm point of $A$ with support $p$.
This construction is a bijection with $\mathcal S(A)$; its inverse uses
the actual support and induced residue-field valuation of a seminorm point.
\end{theorem}
\begin{proof}
Pullback of $\nu$ along $A\to\kappa(p)$ is a multiplicative ultrametric
seminorm. A valuation on a field vanishes precisely at zero, and the
kernel of $A\to\kappa(p)$ is exactly $p$, proving the support statement.
The induced residue-field valuation of a seminorm agrees with the original
values on every $a\in A$, so evaluation after the proposed inverse
recovers the entire seminorm. If two valued points give the same
seminorm, their supports agree, hence their primes agree. Their valuations
on the common residue field agree on all images of $A$; the residue-field
fraction representation makes those valuations equal. Thus evaluation is
injective and the other inverse identity follows as well.
\end{proof}

\begin{theorem}[Canonical evaluation topology]
\label{spectrum:valued-point-topology}
\lean{EnvelopeProperty.Berkovich.ValuedPoint.toFun,
EnvelopeProperty.Berkovich.ValuedPoint.toFun_injective,
EnvelopeProperty.Berkovich.ValuedPoint.isInducing_toFun,
EnvelopeProperty.Berkovich.ValuedPoint.isEmbedding_toFun,
EnvelopeProperty.Berkovich.ValuedPoint.continuous_toFun,
EnvelopeProperty.Berkovich.ValuedPoint.continuous_eval,
EnvelopeProperty.Berkovich.ValuedPoint.continuous_iff,
EnvelopeProperty.Berkovich.ValuedPoint.continuous_toSpectrum,
EnvelopeProperty.Berkovich.ValuedPoint.continuous_ofSpectrum,
EnvelopeProperty.Berkovich.ValuedPoint.spectrumHomeomorph,
EnvelopeProperty.Berkovich.ValuedPoint.spectrumHomeomorph_apply,
EnvelopeProperty.Berkovich.ValuedPoint.spectrumHomeomorph_symm_apply,
EnvelopeProperty.Berkovich.ValuedPoint.continuous_prime}
\leanok
\uses{spectrum:valued-point-equivalence,spectrum:topology,spectrum:support}
Give valued points the coarsest topology making every original-ring
evaluation continuous. The proved point equivalence is a homeomorphism
with $\mathcal S(A)$, and the valued-point space is Hausdorff. Projection
to the actual prime is continuous for the Zariski topology.
\end{theorem}
\begin{proof}
The topology is induced from the map to the product of all evaluation
coordinates. Evaluation is injective by the preceding theorem, so this
is an embedding into a Hausdorff product. A map to the valued-point
space is continuous exactly when all its evaluations are continuous.
Both directions of the equivalence preserve every evaluation, giving
continuity in both directions and the asserted homeomorphism. The prime
projection is the composite of this homeomorphism with the already proved
continuous support map, because the evaluation support is exactly the prime.
\end{proof}

\begin{theorem}[Actual residue-field pullback]
\label{spectrum:valued-point-pullback}
\lean{EnvelopeProperty.Berkovich.ValuedPoint.comap,
EnvelopeProperty.Berkovich.ValuedPoint.comap_prime,
EnvelopeProperty.Berkovich.ValuedPoint.eval_comap,
EnvelopeProperty.Berkovich.ValuedPoint.toSpectrum_comap,
EnvelopeProperty.Berkovich.ValuedPoint.continuous_comap,
EnvelopeProperty.Berkovich.ValuedPoint.comap_id,
EnvelopeProperty.Berkovich.ValuedPoint.comap_comp}
\leanok
\uses{spectrum:valued-point-equivalence,spectrum:valued-point-topology,
spectrum:residue-field-pullback}
A ring homomorphism $\varphi:A\to B$ pulls a valued point $(q,\nu)$
back to the prime $\varphi^{-1}(q)$ with valuation obtained by pulling
$\nu$ back along the actual induced residue-field homomorphism.
Its evaluations satisfy $|a|=|\varphi(a)|$, so this map is continuous
and agrees with seminorm pullback under the homeomorphism. It obeys
identity and composition laws on the whole points.
\end{theorem}
\begin{proof}
The preimage of the prime is prime, and its residue field maps to
$\kappa(q)$ by the residue-field universal property. Pull back the
given valuation along this genuine field map. The residue-field map
sends the image of each $a$ to the image of $\varphi(a)$, proving
the evaluation identity. Continuous evaluation proves continuity,
and equality on all ring elements proves compatibility with the
seminorm equivalence. The identity and composition laws of original
ring pullback, together with injectivity of evaluation, give the laws
for the valued points.
\end{proof}

\begin{theorem}[Valued points with prescribed coefficients]
\label{spectrum:relative-valued-points}
\lean{EnvelopeProperty.Berkovich.ValuedPoint.relativeLocus,
EnvelopeProperty.Berkovich.ValuedPoint.mem_relativeLocus,
EnvelopeProperty.Berkovich.ValuedPoint.isClosed_relativeLocus,
EnvelopeProperty.Berkovich.RelativeValuedPoint,
EnvelopeProperty.Berkovich.RelativeValuedPoint.spectrumHomeomorph,
EnvelopeProperty.Berkovich.RelativeValuedPoint.spectrumHomeomorph_apply_coe,
EnvelopeProperty.Berkovich.RelativeValuedPoint.spectrumHomeomorph_symm_apply_coe,
EnvelopeProperty.Berkovich.RelativeValuedPoint.isClosedEmbedding_subtypeVal,
EnvelopeProperty.Berkovich.RelativeValuedPoint.eval,
EnvelopeProperty.Berkovich.RelativeValuedPoint.eval_algebraMap,
EnvelopeProperty.Berkovich.RelativeValuedPoint.continuous_eval,
EnvelopeProperty.Berkovich.RelativeValuedPoint.prime_base,
EnvelopeProperty.Berkovich.RelativeValuedPoint.comap,
EnvelopeProperty.Berkovich.RelativeValuedPoint.comap_coe,
EnvelopeProperty.Berkovich.RelativeValuedPoint.eval_comap,
EnvelopeProperty.Berkovich.RelativeValuedPoint.spectrumHomeomorph_comap,
EnvelopeProperty.Berkovich.RelativeValuedPoint.continuous_comap,
EnvelopeProperty.Berkovich.RelativeValuedPoint.comap_id,
EnvelopeProperty.Berkovich.RelativeValuedPoint.comap_comp}
\leanok
\uses{spectrum:valued-point-topology,spectrum:valued-point-pullback,spectrum:relative,
spectrum:relative-support}
For any commutative base ring $R$ carrying a prescribed real seminorm
$v$, valued points of an $R$-algebra $A$ whose evaluations on all
$R$ coefficients equal $v$ form a closed subspace. The valued-point
homeomorphism restricts to a homeomorphism with $\mathcal S_v(A)$. Supports
contract to the support of $v$. Actual residue-field pullback along an
$R$-algebra map is continuous, preserves the prescribed coefficients, and
agrees with relative seminorm pullback, including identity and composition laws.
\end{theorem}
\begin{proof}
Prescribing each coefficient evaluation is a closed equality condition,
so their intersection is closed. The absolute homeomorphism preserves
every coefficient evaluation, hence identifies the two prescribed-base
subspaces exactly. Restricting the genuine homeomorphism gives the
relative homeomorphism. A coefficient has zero value exactly when its image
lies in the valued-point prime, so this prime contracts to the support of $v$.
An algebra map commutes with coefficient inclusion, and its genuine
residue-field pullback therefore preserves every prescribed coefficient
value. The absolute evaluation identity proves relative naturality, continuity,
and both functorial laws. No existence over an arbitrary prime or base
seminorm is assumed in this construction.
\end{proof}
